% !TeX root = ../../Book1.tex
% 纲要标题：### 13.5 叶状结构的整体问题与不可积模型
% 纲要位置：几何学计划纲要.md，第 635 行。
% 正文状态：初稿；尚未正式审核。

% 纲要内容开始
% 
% * 环面上的有理和无理斜率叶
% * 横截面、返回映射及沿叶映射芽
% * 非 Hausdorff 叶空间
% * 标准接触分布的不可积性
% 
% ---
% 
% 纲要内容结束

\section{叶状结构的整体问题与不可积模型}
\label{g1:sec:1305}

Frobenius 定理给每个点附近一叠坐标切片，最大叶定理则说明怎样沿这些切片继续前进。继续以后会发生什么，仍须另外研究：一条叶可能闭合，也可能反复进入同一个坐标盒，甚至在整个流形中稠密；沿叶绕一圈后，邻近叶的横向位置也可能改变。本节用几个能够直接计算的模型分辨这些现象。全节的流形均无边界，分布均为常秩光滑分布。

\subsection{环面上的有理和无理斜率叶}
\label{g1:subsec:1305:01}

把二维环面写为 \(T^2=\mathbb R^2/\mathbb Z^2\)，即分别把两个坐标模 \(1\) 识别。局部坐标来自小矩形的商像，换坐标只是整数平移，所以常向量场
\[
	X_\alpha=\partial_x+\alpha\partial_y,\qquad \alpha\in\mathbb R,
\]
下降为环面上的处处非零光滑向量场。它的流为
\[
	\Phi_t([x,y])=[x+t,y+\alpha t].
\]
这给出一个处处秩一的分布 \(D_\alpha=\mathbb RX_\alpha\)。秩一分布总对合，因为 \([fX_\alpha,gX_\alpha]=(fX_\alpha g-gX_\alpha f)X_\alpha\)；由\cref{g1:thm:maximal-connected-leaf}，其叶就是这些流轨道，带上各自的积分流形拓扑。

\begin{proposition}[环面上的线性叶]
	\label{g1:prop:linear-torus-leaves}
	设 \(T^2=\mathbb R^2/\mathbb Z^2\)，\(\alpha\in\mathbb R\)，\(X_\alpha=\partial_x+\alpha\partial_y\)，\(D_\alpha=\mathbb R X_\alpha\)，其流为 \(\Phi_t([x,y])=[x+t,y+\alpha t]\)。若 \(\alpha=p/q\)，其中 \(p\in\mathbb Z\)、\(q\in\mathbb Z_{\geq1}\) 且 \(\gcd(p,q)=1\)，则每条 \(D_\alpha\) 叶为嵌入圆周，最小正流周期为 \(q\)。若 \(\alpha\) 无理，则每条叶为 \(\mathbb R\) 的单射浸入像，在整个环面中稠密，且不是嵌入子流形。
\end{proposition}
\begin{proof}
	轨道在时刻 \(t\) 返回初始点，当且仅当 \(t\in\mathbb Z\) 且 \(\alpha t\in\mathbb Z\)。有理情形中这等价于 \(t\in q\mathbb Z\)。于是轨道参数化下降为单射浸入 \(\mathbb R/q\mathbb Z\to T^2\)。定义域紧而目标 Hausdorff，所以它到像为同胚，结合浸入性即为光滑嵌入。

	无理情形没有非零返回时间，故轨道参数化 \(\mathbb R\to T^2\) 单射，其微分是处处非零的 \(X_\alpha\)，因此是浸入。先证明整数倍 \(\{n\alpha\bmod1:n\in\mathbb Z\}\) 在圆周稠密。将 \(0,\alpha,\ldots,N\alpha\) 的小数部分放入 \(N\) 个等长区间，至少两点落在同一段。两者之差取适当正负号，就得到某个非零整数 \(k\) 及整数 \(m\)，使
	\[
		0<\delta=k\alpha-m<\frac1N.
	\]
	各数 \(0,\delta,\ldots,\lfloor1/\delta\rfloor\delta\) 都来自整数倍 \(\alpha\) 模 \(1\)，并且与 \([0,1]\) 中任一点的距离不超过 \(\delta\)。令 \(N\) 任意大，便得稠密性。对环面中任意小开矩形，先固定其中一个第一坐标 \(u\)。轨道在时刻 \(t=u-x_0+n\) 的第一坐标均为 \(u\) 模 \(1\)，第二坐标为
	\[
		y_0+\alpha(u-x_0)+n\alpha\pmod1,
	\]
	这些值稠密，所以某个时刻落入给定矩形。这证明整条叶稠密。

	还须辨认叶的拓扑。由同一逼近论证，可以选 \(|n_j|\to\infty\)，使 \(n_j\alpha\bmod1\to0\)；若这些整数的绝对值有界，则其到整数的非零距离有正的最小值，与任意精度逼近矛盾。于是 \(\Phi_{n_j}(p)\to p\) 在环面中成立，但参数 \(n_j\) 不趋于零。故参数化的逆在子空间像上不连续，它不是嵌入。另一方面，参数化到最大叶是双射局部微分同胚，因为在每个流盒中都是时间坐标的平移，故与叶自身的光滑结构微分同胚。叶本身为 \(\mathbb R\)，稠密性发生在它到环面的浸入像上。
\end{proof}

% !TeX root = ../Book1.tex
\begin{figure}[htbp]
	\centering
	\begin{tikzpicture}[x=3.65cm,y=3.65cm,>=Stealth,font=\small]
		\begin{scope}[shift={({0},0)}]
			\fill[blue!3] (0,0) rectangle (1,1);
			\draw[gray] (0,0) rectangle (1,1);
			\draw[blue!70!black,semithick] (0.0000,0.1400) -- (1.0000,0.6400);
			\draw[blue!70!black,semithick] (0.0000,0.6400) -- (0.7200,1.0000);
			\draw[blue!70!black,semithick] (0.7200,0.0000) -- (1.0000,0.1400);
			\fill[orange!85!black] (0,.14) circle[radius=.016];
			\node[below left] at (0,0) {$0$};
			\node[below] at (1,0) {$1$};
			\node[left] at (0,1) {$1$};
			\node at (.5,1.13) {有理斜率 $\alpha=1/2$};
			\node at (.5,-.18) {一周闭合轨迹};
		\end{scope}
		\begin{scope}[shift={({1.55},0)}]
			\fill[blue!3] (0,0) rectangle (1,1);
			\draw[gray] (0,0) rectangle (1,1);
			\draw[blue!70!black,semithick] (0.0000,0.1400) -- (0.6081,1.0000);
			\draw[blue!70!black,semithick] (0.6081,0.0000) -- (1.0000,0.5542);
			\draw[blue!70!black,semithick] (0.0000,0.5542) -- (0.3152,1.0000);
			\draw[blue!70!black,semithick] (0.3152,0.0000) -- (1.0000,0.9684);
			\draw[blue!70!black,semithick] (0.0000,0.9684) -- (0.0223,1.0000);
			\draw[blue!70!black,semithick] (0.0223,0.0000) -- (0.7294,1.0000);
			\draw[blue!70!black,semithick] (0.7294,0.0000) -- (1.0000,0.3826);
			\draw[blue!70!black,semithick] (0.0000,0.3826) -- (0.4365,1.0000);
			\draw[blue!70!black,semithick] (0.4365,0.0000) -- (1.0000,0.7969);
			\draw[blue!70!black,semithick] (0.0000,0.7969) -- (0.1436,1.0000);
			\draw[blue!70!black,semithick] (0.1436,0.0000) -- (0.8508,1.0000);
			\draw[blue!70!black,semithick] (0.8508,0.0000) -- (1.0000,0.2111);
			\draw[blue!70!black,semithick] (0.0000,0.2111) -- (0.5579,1.0000);
			\draw[blue!70!black,semithick] (0.5579,0.0000) -- (1.0000,0.6253);
			\draw[blue!70!black,semithick] (0.0000,0.6253) -- (0.2650,1.0000);
			\draw[blue!70!black,semithick] (0.2650,0.0000) -- (0.9721,1.0000);
			\draw[blue!70!black,semithick] (0.9721,0.0000) -- (1.0000,0.0395);
			\fill[orange!85!black] (0,.14) circle[radius=.016];
			\node[below left] at (0,0) {$0$};
			\node[below] at (1,0) {$1$};
			\node[left] at (0,1) {$1$};
			\node at (.5,1.13) {无理斜率 $\alpha=\sqrt2$};
			\node at (.5,-.18) {有限时间轨迹};
		\end{scope}
	\end{tikzpicture}
	\caption[环面上线性流的两种整体行为]{正方形的相对边分别识别，线段在右边离开后从左边进入，在上边离开后从下边进入。左图画出一次闭合的轨迹；右图只画出有限时间内的轨迹，稠密性由正文的整数逼近证明。两图中的每条局部线段都只有一维。}
	\label{g1:fig:torus-leaves}
\end{figure}



局部流盒看不出斜率是否有理：两种情形都能把向量场写成 \(\partial_s\)。区别来自绕过环面后坐标的整数识别。无理叶也没有在一个流盒中占据整块二维区域；它在盒内出现许多互不相连的局部线段，这些线段的并可以稠密，而每一段仍只有一维。

\subsection{横截面、返回映射及沿叶映射芽}
\label{g1:subsec:1305:02}

一条轨道离开小流盒后怎样回来，可以通过它与横截面的交点来记录。附近轨道也给出相应交点，因而沿流运动在横截面之间诱导局部映射。

\begin{definition}
	\label{g1:def:foliation-transversal}
	设 \(D\subset \mathrm{T}M\) 是秩 \(r\) 的可积分布，\(\dim M=n\)。一个 \((n-r)\) 维局部嵌入子流形 \(S\) 若满足
	\[
		\mathrm{T}_pM=D_p\oplus \mathrm{T}_pS\qquad(p\in S),
	\]
	称为这族叶的\term[hengjiemian]{横截面}[transversal]。对于处处非零向量场产生的秩一分布，这表示向量场不切于 \(S\)。
\end{definition}

在适应坐标 \((x,y)\) 中，叶片由 \(y\) 固定给出，横截面可以取 \(x=0\)。由互补条件，横向投影 \((x,y)\mapsto y\) 限制到任意横截面时，其微分可逆；\cref{g1:thm:inverse-function}说明它局部为微分同胚。所以在一个足够小的叶状坐标盒中，附近叶片与横截面的交点可用横向坐标唯一标记。

\begin{proposition}[横截面间的局部返回]
	\label{g1:prop:local-return-map}
	设 \(X\) 为无边界光滑流形 \(M\) 上的光滑向量场，\(\Phi\) 为它的最大局部流，\(\Phi_T(p)=q\)，其中 \(T>0\)，且 \(X\) 沿这段轨道不为零。在 \(p,q\) 附近分别取横截面 \(S,S'\)。存在 \(S\) 中 \(p\) 的邻域 \(V\) 及光滑函数 \(\tau:V\to\mathbb R\)，满足 \(\tau(p)=T\)，使
	\[
		P(z)=\Phi_{\tau(z)}(z)\in S'
	\]
	为到 \(S'\) 中 \(q\) 的一个邻域的微分同胚。这里 \(\tau(z)\) 是接近指定时间 \(T\) 的交会时刻，未必是第一次正返回时间。
\end{proposition}
\begin{proof}
	在 \(q\) 附近将 \(S'\) 写成 \(h=0\)，其中 \(\mathrm{d}h_q\neq0\)。横截性等价于 \(\mathrm{d}h_q(X_q)\neq0\)。\cref{g1:thm:maximal-local-flow}保证 \((t,z)\mapsto\Phi_t(z)\) 在 \((T,p)\) 附近光滑，于是方程
	\[
		H(t,z)=h(\Phi_t(z))=0
	\]
	满足 \(H(T,p)=0\)、\(\partial_tH(T,p)=\mathrm{d}h_q(X_q)\neq0\)。参数隐函数定理给出唯一接近 \(T\) 的光滑解 \(t=\tau(z)\)。这里可在 \(S\) 的局部坐标中直接应用\cref{g1:thm:implicit-function-parameters}。

	沿同一轨道反向运动，用 \(S\) 的局部定义函数作相同构造，就得到 \(S'\) 中 \(q\) 附近到 \(S\) 的光滑交会映射 \(Q\)，时间接近 \(-T\)。适当缩小邻域后，\(Q(P(z))=z\)：原来的 \(z\) 已是接近 \(-T\) 时刻的交点，而隐函数定理保证此交点唯一。同理 \(P(Q(w))=w\)。故两映射互逆。
\end{proof}

若 \(q=p\) 且 \(S'=S\)，就得到闭轨道附近的\term[fanhuiyingshe]{返回映射}[return map]。一般轨道可能不再遇到选定截面，也可能在所选 \(T\) 之前遇到截面；存在一个局部交会映射，不等于全体点具有共同定义的首次返回映射。在环面线性流中，截面 \(S=\{x=0\}\) 则确实是整体横截面，所有点的首次正返回时间都是 \(1\)，返回映射为
\[
	P([y])=[y+\alpha].
\]
于是有理斜率给出有限次迭代后的返回，无理斜率给出无周期的圆周旋转。

\ProseParagraphBreak

高维叶没有统一的“沿叶前进时间”，但仍可指定一条叶内道路，并看附近叶怎样跟随它穿过一串坐标盒。设 \(\gamma:[0,1]\to L\) 是一条叶内道路，端点为 \(p,q\)，在两端固定横截面 \(S_p,S_q\)。取有限分割，使每小段道路都落在一张叶状坐标盒内。沿每一段保持横向坐标不变，抵达下一张小横截面；在交叠处使用光滑换坐标，再作下一段。有限次复合得到 \(p\) 附近到 \(q\) 附近的局部微分同胚。

\begin{definition}
	\label{g1:def:leaf-holonomy}
	设 \(D\) 为无边界光滑流形 \(M\) 上的常秩光滑可积分布，\(L\) 为 \(D\) 的一条最大连通积分叶，\(\gamma:[0,1]\to L\) 为叶内道路，\(p=\gamma(0)\)、\(q=\gamma(1)\)，并在两端固定局部横截面 \(S_p,S_q\)。沿覆盖道路的有限个适应坐标盒，在每一张盒中按相同横向坐标作叶片对应，再依次复合，得到 \(S_p\) 中 \(p\) 附近到 \(S_q\) 中 \(q\) 附近的局部微分同胚。两个这样的局部映射若在 \(p\) 的某个更小邻域相同，就称它们在 \(p\) 有相同的映射芽。上述复合在 \(p\) 的芽记为 \(h_\gamma\)，称为\term[yeholonomy]{沿叶道路的横截面映射芽}[leaf holonomy]。当 \(\gamma\) 为基于 \(p\) 的环路时，在两端采用同一横截面，所得 \(h_\gamma\) 是固定 \(p\) 的横截面返回映射芽；所有这样的芽在复合下组成该叶在 \(p\) 的横截面返回映射芽群。
\end{definition}

采用芽的原因是，改变坐标盒大小会改变代表映射的定义域，却不应改变靠近初始点的横向效应。这里移动的是邻近叶的交点，并没有选择联络，也没有沿曲线搬运切向量；它与\chapref{g1:ch:09}的环路平行移动变换是不同的构造。

\begin{proposition}[沿叶横截面映射芽的相容性]
	\label{g1:prop:leaf-holonomy-homotopy}
	设 \(D\) 为无边界光滑流形上的常秩光滑可积分布，固定道路两端的局部横截面。沿叶道路的横截面映射芽不依赖所选坐标盒与分割，只依赖道路在该叶中相对端点的同伦类。若 \(\gamma*\eta\) 表示先走 \(\gamma\) 后走 \(\eta\)，则
	\[
		h_{\gamma*\eta}=h_\eta\circ h_\gamma,\qquad
		h_{\gamma^{-1}}=h_\gamma^{-1}.
	\]
	改变基点处横截面，用两截面间沿叶片的局部对应识别后，环路产生的芽按共轭变换。
\end{proposition}
\begin{proof}
	先看一张小坐标盒。在适应坐标中，叶片是 \(y=\text{常数}\)，各横截面经投影到 \(y\) 坐标都是局部微分同胚。因此两横截面间的对应是“投影到 \(y\)，再用另一截面的投影逆”，只由两截面决定。插入一张中间截面时，两个投影互相消去，复合不变。

	在两张适应图交叠处，第二张图的横向坐标对第一张图的叶方向导数为零，因为两图都将这些方向识别为 \(D\)。缩小到交叠内的乘积小盒后，它因此只依赖第一张图的横向坐标；这个横向变化的微分可逆，因为整个换坐标微分可逆且有相应的块三角形式。故同一小段道路用任一张图计算，得到相同的横截面对应芽。

	对一条固定道路，将两种有限分割取公共加细，再把每段缩小到上述交叠小盒中。插入中间截面和改变坐标均不改变对应芽，有限次复合仍相同。这证明了定义与选择无关。拼接和逆向的公式也由依次复合与逐段求逆得到；所有比较只需在有限个代表映射的共同小定义域内进行。

	再设 \(H:[0,1]^2\to L\) 是固定端点的叶内同伦。这里连续性使用叶自身的拓扑。每个点都有叶中开邻域，其像落在一张叶片内；这些邻域的原像覆盖紧参数方形。有限子覆盖和 Lebesgue 数给出足够细的网格，使每个小格的像都落在同一张叶片内。为说明所用紧性步骤，可对有限开覆盖选每点一个仍留在某成员中的小球，取有限个半径减半的球覆盖方形；这些半径的正最小值就给出一个可用的细分尺度。

	每个小格边界的返回映射芽 是恒等芽，因为在同一张叶片内，沿边的所有横截面对应都由同一个横向坐标决定，绕一圈后投影及其逆全部抵消。把网格逐格拼起来，公共边的两个反向对应也互相抵消，得到方形外边界的返回映射芽为恒等芽。两条竖边因端点固定而给出恒等芽，所以上下两条道路的横截面映射芽相同。最后，若横截面通过局部对应 \(c\) 改变，沿环路前后分别插入 \(c^{-1}\) 与 \(c\)，所得芽为 \(c\circ h_\gamma\circ c^{-1}\)，即共轭。
\end{proof}

\begin{example}[收缩的横向返回]
	\label{g1:exm:contracting-leaf-holonomy}
	在圆柱 \((\mathbb R/\mathbb Z)\times\mathbb R\) 上取
	\[
		X=\partial_\theta-a r\partial_r,\qquad a>0.
	\]
	其整体流为 \(\Phi_t(\theta,r)=(\theta+t,e^{-at}r)\)。中间圆周 \(r=0\) 是闭叶，截面 \(\theta=0\) 的返回映射为 \(P(r)=e^{-a}r\)。沿闭叶正向一圈的横截面返回映射芽因而是这个收缩映射在零点的芽；绕 \(k\) 圈得到 \(r\mapsto e^{-ak}r\)。它不是恒等芽，反映的是邻近叶的横向收缩。

	相比之下，无理斜率环面的每条叶自身是 \(\mathbb R\)，所有叶内环路均可缩，所以由上一命题，横截面返回映射芽群平凡。截面上的无理旋转返回把一个交点送到另一个交点，并非沿该叶环路返回同一个叶内基点。这说明返回运动复杂与沿环路的横截面返回映射芽非平凡不是同一件事。
\end{example}

\subsection{非 Hausdorff 叶空间}
\label{g1:subsec:1305:03}

把每条叶压成一点后，叶的聚集行为就转化为商空间的邻域关系。开集能否将不同叶分开，将决定这个商能否作为通常意义下的流形继续研究。

\begin{definition}
	\label{g1:def:leaf-space}
	设 \(\mathcal F\) 为光滑流形 \(M\) 上的光滑叶状结构。将同一叶中的所有点识别，得到的集合记为 \(M/\mathcal F\)，称为\term[yekongjian]{叶空间}[leaf space]，以 \(q:M\to M/\mathcal F\) 表示将每点送到其所在叶的映射，并赋予商拓扑。即 \(V\subset M/\mathcal F\) 开，当且仅当\(q^{-1}(V)\) 在 \(M\) 中开；这样的逆像是若干整条叶的并，也称饱和开集。
\end{definition}

这个拓扑描述的是能否用开区域区分整条叶，与上一节每条叶自身的流形拓扑不同。在有理斜率环面上，叶空间确实是圆周。令 \(\alpha=p/q\) 为既约分数，定义
\[
	f:T^2\longrightarrow\mathbb R/\mathbb Z,\qquad f([x,y])=[qy-px].
\]
它良定义、满射，且沿 \(X_\alpha\) 的导数为零。若两点的 \(f\) 值相同，取实代表后有 \(q\Delta y-p\Delta x=k\in\mathbb Z\)。互素性给出整数 \(m,n\) 满足 \(qn-pm=k\)，从而 \(q(\Delta y-n)=p(\Delta x-m)\)。这恰好说明两点在同一条斜率 \(p/q\) 的轨道上。因此纤维就是叶，\(f\) 诱导叶空间到圆周的连续双射。叶空间是紧环面的连续像，故紧；圆周 Hausdorff，这个连续双射就是同胚。

\begin{proposition}[无理斜率的叶空间]
	\label{g1:prop:irrational-leaf-quotient}
	设 \(T^2=\mathbb R^2/\mathbb Z^2\)，\(\alpha\in\mathbb R\setminus\mathbb Q\)，\(\mathcal F_\alpha\) 为分布 \(\operatorname{span}\{\partial_x+\alpha\partial_y\}\) 的积分叶所成的叶状结构。叶空间 \(T^2/\mathcal F_\alpha\)不止一个点，但只有空集和全空间两个开集。特别地，它不是 Hausdorff 空间。
\end{proposition}
\begin{proof}
	每条叶与横截面 \(x=0\) 的交点集为一个整数旋转轨道，至多可数；圆周上的点不可数，所以不止一条叶。若 \(U\subset T^2\) 是非空饱和开集，每条叶都稠密，因此每条叶都与 \(U\) 相交。饱和性又要求整条相交的叶包含于 \(U\)，所以 \(U=T^2\)。按商拓扑的定义，这恰好给出所述两个开集。不同叶所对应的两点没有不交邻域，故非 Hausdorff。
\end{proof}

叶空间的分离失败并不只发生在稠密叶中。下面所有叶都是闭的嵌入子流形，叶空间仍然无法分离某两个点。

\begin{example}[分裂的零层]
	\label{g1:exm:punctured-plane-leaf-space}
	令 \(M=\mathbb R^2\setminus\{(0,0)\}\)，\(D=\mathbb R\partial_x\)。当 \(y\neq0\) 时，直线 \(\{(x,y):x\in\mathbb R\}\) 是一条叶；当 \(y=0\) 时，缺失的原点将该直线分成两条叶
	\[
		L_- =\{(x,0):x<0\},\qquad L_+=\{(x,0):x>0\}.
	\]
	这些叶都在 \(M\) 中闭，且以其通常参数为嵌入子流形。例如 \(L_+\) 的环境闭包可能增加的唯一有限点是原点，而原点已被删除。

	但 \(L_-\) 与 \(L_+\) 无法被不交的饱和开集分开。若饱和开集 \(U_-\) 包含 \(L_-\)，它包含 \((-1,0)\) 的一个小开矩形，因此对所有充分小的 \(y\neq0\)，\(U_-\) 都与高度 \(y\) 的水平叶相交，进而包含整条该叶。包含 \(L_+\) 的饱和开集 \(U_+\) 同样包含所有充分小非零高度的整条水平叶。所以 \(U_-\cap U_+\neq\varnothing\)。叶空间因此非 Hausdorff，尽管每个叶类是闭集、商空间的每个单点也是闭集。
\end{example}

局部横向坐标仍然存在；问题是绕过缺口继续这些坐标时，零高度对应了两条不同叶，而邻近非零高度各只对应一条叶。构造一个整体流形作为叶空间，需要超出局部可积性的额外条件。

\subsection{标准接触分布的不可积性}
\label{g1:subsec:1305:04}

前面的例子都有积分叶，只是叶的整体组织方式不同。现在看另一种障碍：一族光滑平面甚至不能在任何一点附近拼成与它们相切的曲面。在 \(\mathbb R^3\) 中定义
\[
	\alpha=\mathrm{d}z-x\,\mathrm{d}y,\qquad
	D=\ker\alpha
	=\operatorname{span}\{A,B\},\qquad
	A=\partial_x,\quad B=\partial_y+x\partial_z.
\]
因为 \(\alpha(\partial_z)=1\)，这确为常秩二的光滑分布。\(A\) 允许沿 \(x\) 方向移动，\(B\) 允许沿 \(y\) 方向移动，但后者同时按当前位置 \(x\) 改变高度。

\begin{proposition}[标准接触平面的不可积性]
	\label{g1:prop:contact-model-nonintegrable}
	在 \(\mathbb R^3\) 的标准坐标 \((x,y,z)\) 中，令 \(\alpha=\mathrm{d}z-x\,\mathrm{d}y\)，\(A=\partial_x\)，\(B=\partial_y+x\partial_z\)，\(D=\ker\alpha=\operatorname{span}\{A,B\}\)。分布 \(D\) 在每一点都不对合，因而不存在任何二维积分流形。具体地，\([A,B]=\partial_z\notin D\)。
\end{proposition}
\begin{proof}
	由括号坐标公式，只有 \(B\) 的 \(\partial_z\) 分量 \(x\) 沿 \(A=\partial_x\) 求导产生非零项，故 \([A,B]=\partial_z\)。而 \(\alpha(\partial_z)=1\)，所以该向量不属于核。若存在二维积分流形，则其上任一点都必须满足\cref{g1:prop:involutivity-necessary}的必要条件，与此矛盾。

	还可直接看出曲面方程的矛盾。若有积分曲面，投影 \((x,y,z)\mapsto(x,y)\) 在其切空间上可逆，因为它把 \(A,B\) 送为平面坐标基。逆函数定理遂将曲面局部写成 \(z=f(x,y)\)。切向条件 \(\alpha|_{\mathrm{T}S}=0\) 要求 \(f_x=0\)、\(f_y=x\)，从而 \(f_{xy}=0\) 而 \(f_{yx}=1\)，违反光滑函数混合偏导相等。这把括号障碍解释成了相切方程的不相容。
\end{proof}

为什么这个模型称为接触分布？在三维中，可以直接用括号给出不依赖外微分的定义。设 \(D\) 为秩二分布，对 \(u,v\in D_p\)，选局部 \(D\) 截面 \(U,V\) 延拓它们，令
\[
	\mathcal B_p(u,v)=[U,V]_p\pmod{D_p}\in \mathrm{T}_pM/D_p.
\]
这个值不依赖延拓。若 \(U\) 改变为 \(U+\sum_i f_iE_i\)，其中 \(E_i\) 为 \(D\) 的局部标架且 \(f_i(p)=0\)，则括号的变化在 \(p\) 为 \(-\sum_iV(f_i)(p)E_i(p)\)，属于 \(D_p\)；改变 \(V\) 同理。因此 \(\mathcal B_p\) 是良定义的交替双线性映射。若它在每点非零，称 \(D\) 为三维流形上的\term[jiechufenbu]{接触分布}[contact distribution]。因 \(D_p\) 为二维、商空间为一维，非零也等价于这个交替配对非退化。上述模型中 \(\mathcal B_p(A_p,B_p)=[\partial_z]\neq0\)，所以处处满足定义。

\ProseParagraphBreak

接触条件比“某处不对合”更强，它要求在每点都有这种非退化的横向括号效应。反过来，对合条件正是 \(\mathcal B_p=0\) 对所有 \(p\) 成立。这一比较把可积与接触两端放在同一个逐点量中：前者的允许方向在交换后不产生新方向，后者的两个允许方向交换后能产生缺少的第三个方向。\chapref{g1:ch:14}建立外微分后，三维接触条件还可写成 \(\alpha\wedge\mathrm{d}\alpha\neq0\)。

% 纲要条目排为编号小节标题，并进入目录与书签；说明段仍保留于纲要注释中。
